%%%%%%%%%%  scan idx 183  |  folio 170  %%%%%%%%%%
\SGAPrintPage{170}%
$F\,|\,S_{f_{ij}}$ is
generated by finitely many sections $h_{ijk}$ over $S_{f_{ij}}$. There exists (EGA I 9.3.1 and IV 1.7.5) an integer $q_i\ge0$ such that $h_{ijk}\otimes f_{ij}^{\otimes q_i}$ extends to a section of $F\otimes L_i^{\otimes q_im_i}$ over $S$. In other words, $F\otimes L_i^{\otimes q_im_i}$ is generated over $S_i$ by finitely many of its sections over $S$. Proceeding as in the proof of (EGA II 4.5.5), we deduce that there exists an integer $n_i$, which we may take to be $>0$, such that, for every $n\ge n_i$, $F\otimes L_i^{\otimes n}$ is generated over $S_i$ by finitely many of its sections over $S$. In particular, there exist an integer $r_i\ge0$ and a morphism
\[
\underline{O}_S^{r_i}\otimes L_i^{\otimes-n_i}\longrightarrow F
\]
which induces an epimorphism over $S_i$. The implication (i$'$)$\Rightarrow$(ii) follows immediately.

(ii)$\Rightarrow$(iii): Since every quasi-coherent sheaf on $S$ is an inductive limit of its quasi-coherent subsheaves of finite type (EGA I 9.4.9 and IV 1.7.7), we may restrict ourselves to proving (iii) for $F$ of finite type. There then exist, for $i\in I$, morphisms
\[
u_i:\underline{O}_S^{r_i}\otimes L_i^{\otimes-n_i}\longrightarrow F
\]
such that $\sum u_i$ is an epimorphism. Now we have
\[
\begin{aligned}
\operatorname{Im}(u_i)&=\operatorname{Im}\bigl(u_i\otimes\operatorname{Id}(L_i^{\otimes n_i})\otimes\operatorname{Id}(L_i^{\otimes-n_i})\bigr)\\
&=\operatorname{Im}\bigl(u_i\otimes\operatorname{Id}(L_i^{\otimes n_i})\bigr)\otimes L_i^{\otimes-n_i}
\end{aligned}
\]
But $\operatorname{Im}(u_i\otimes\operatorname{Id}(L_i^{\otimes n_i}))$ is a subsheaf of $F_{(i,n_i)}$, whence the assertion.

(iii)$\Rightarrow$(iii$'$): This is trivial.

(iii$'$)$\Rightarrow$(i): Let $x\in S$, $U$ be an open subset containing $x$, and $F$ a
%%%%%%%%%%  scan idx 184  |  folio 171  %%%%%%%%%%
\SGAPrintPage{171}%
quasi-coherent Ideal
defining $S-U$. By (iii$'$), there exist $(i,n)\in I\times\mathbb N^+$ and a section $f$ of $F\otimes L_i^{\otimes n}$ over $S$ such that $f(x)\ne0$. We have $x\in S_f\subset U$, which proves (i).

\sgaStableTarget{sga6.e02.def.2-2-4}\textbf{Definition 2.2.4.} \uline{Let} $S$ \uline{be a quasi-compact and quasi-separated scheme, and let} $(L_i)_{i\in I}$ \uline{be a family of invertible sheaves on} $S$\uline{. We say that the} $L_i$ \uline{form an} ample family\SGAAuthorFootnote{raisedstar}{This notion, in form (i$'$), was introduced by M.~BORELLI and S.~L.~KLEIMAN; the terminology \textquotedblleft ample family\textquotedblright{} is proposed by the editor.} \uline{when the equivalent conditions of} (2.2.3) \uline{hold. (In particular, a family consisting of a single sheaf} $L$ \uline{is ample precisely when} $L$ \uline{is ample in the usual sense} (EGA II 4.5.3)).

\sgaStableTarget{sga6.e02.def.2-2-5}\textbf{Definition 2.2.5.} \uline{A scheme} $S$ \uline{is said to be} divisorial \uline{when} $S$ \uline{is quasi-compact and quasi-separated and there exists an ample family} (2.2.4) \uline{of invertible} $\underline{O}_S$\uline{-Modules.}

\sgaStableTarget{sga6.e02.lem.2-2-6}\textbf{Lemma 2.2.6.} \uline{Let} $X$ \uline{be a locally noetherian scheme. Let} $U$ \uline{be an open subset of} $X$ \uline{such that the inclusion} $U\to X$ \uline{is affine (for example} $U$ \uline{affine and} $X$ \uline{separated). Then} $X-U$ \uline{is purely of codimension} $0$ \uline{or} $1$ \uline{(i.e.\ at every maximal point} $x$ \uline{of} $X-U$ \uline{we have} $\dim(\underline{O}_{X,x})=0$ or $1$\uline{), hence purely of codimension} $1$ \uline{when} $U$ \uline{contains the maximal points of} $X$.

\uline{Proof.} Put $Y=X-U$. Let $x$ be a maximal point of $Y$, and make the base change $X'=\operatorname{Spec}(\underline{O}_{X,x})\to X$: we obtain a cartesian diagram:
\[
\begin{tikzcd}
U' \arrow[r] \arrow[d] & X' \arrow[d] & Y' \arrow[l] \arrow[d] \\
U \arrow[r] & X & Y \arrow[l]
\end{tikzcd},
\]
where $Y'=X'-U'$. Since $x$ is a maximal point of $Y$, $Y'$ is precisely the closed point of $X'$. On the other hand, since $X'$ is affine, $U'$ is an affine open subset of $X'$. Let us show that the dimension $n$ of $X'$ is equal to $0$ or $1$. Indeed, we know (SGA 2 V 5.7) that $H^n_{Y'}(X',\underline{O}_{X'})\ne0$. But, for $n\ge2$, we have
\[
H^n_{Y'}(X',\underline{O}_{X'})\simeq H^{n-1}(U',\underline{O}_{U'}),
\]
whereas
\[
H^{n-1}(U',\underline{O}_{U'})=0,
\]
%%%%%%%%%%  scan idx 185  |  folio 172  %%%%%%%%%%
\SGAPrintPage{172}%
since $U'$ is affine and $n-1\ge1$, whence a contradiction.

\sgaStableTarget{sga6.e02.rem.2-2-6-1}\textbf{Remark 2.2.6.1.} Without the separation hypothesis on $X$, (2.2.6) for an affine open subset $U$ of $X$ would be false, as shown by the example of the affine plane with the origin doubled (EGA I 5.5.11).

\sgaStableTarget{sga6.e02.prop.2-2-7}\textbf{Proposition 2.2.7.} \uline{Let} $S$ \uline{be a noetherian, separated, locally factorial scheme (i.e.\ one whose local rings are factorial). Then} $S$ \uline{is divisorial} (2.2.5).

In particular (EGA IV 21.11.1):

\sgaStableTarget{sga6.e02.cor.2-2-7-1}\textbf{Corollary 2.2.7.1.} \uline{Every noetherian regular separated scheme is divisorial.}

\uline{Proof} of 2.2.7. We may assume that $S$ is integral. Let $(S_i)_{i\in I}$ be a finite covering of $S$ by nonempty affine open subsets. By (2.2.6), $S-S_i$ is purely of codimension $1$, hence, by (EGA IV 21.7.4), is the support of a positive divisor. In other words, for each $i\in I$, there exist an invertible sheaf $L_i$ and a section $f_i\in\Gamma(S,L_i)$ such that $S_i=S_{f_i}$. The family $(L_i)$ is then ample (2.2.4) by (2.2.3 (i$'$)). This completes the proof.

\sgaStableTarget{sga6.e02.rem.2-2-7-2}\textbf{Remark 2.2.7.2.} As in (2.2.6), the separation hypothesis on $S$ is essential. Indeed, again take for $S$ the affine plane with the origin doubled. Recall that $S$ is obtained by gluing two copies of the affine plane $\underline{\underline{E}}^2_k$ ($k$, a field) by the identity on the complement $U$ of the origin $O$; we therefore have a canonical projection
\[
p:S\longrightarrow\underline{\underline{E}}^2_k,
\]
which induces an isomorphism over $U$, while $p^{-1}(O)$ consists of two points $O_1$ and $O_2$. An invertible sheaf on $S$ is obtained by gluing two invertible sheaves on $\underline{\underline{E}}^2_k$ by an isomorphism over $U$; since every invertible sheaf on $\underline{\underline{E}}^2_k$ is trivial, an invertible sheaf on $S$ is therefore defined by an element $g\in\Gamma(U,\underline{O}_U^*)$; but $g$ is necessarily a constant, since the restriction $\Gamma(\underline{\underline{E}}^2_k,\underline{O}_{\underline{\underline{E}}^2_k})\to\Gamma(U,\underline{O}_U)$ is an isomorphism, $O$ being of depth $2$ in the plane; thus every
%%%%%%%%%%  scan idx 186  |  folio 173  %%%%%%%%%%
\SGAPrintPage{173}%
invertible sheaf on $S$ is trivial. On the other hand, it is immediate that $p$ induces an isomorphism
\[
p^*:\Gamma(\mathbb E^2_k,\underline{O}_{\mathbb E^2_k})\longrightarrow\Gamma(S,\underline{O}_S).
\]
Consequently, every global section of $\underline{O}_S$ which vanishes at $O_1$ also vanishes at $O_2$, so the $S_f$, for $f\in\Gamma(S,\underline{O}_S)$, do not form a basis for the topology of $S$.

\sgaStableTarget{sga6.e02.thm.2-2-8}\textbf{Lemma 2.2.8.} \uline{Let} $S$ \uline{be a divisorial scheme} (2.2.5), $(L_i)_{i\in I}$ \uline{an ample family of invertible} $\underline{O}_S$\uline{-Modules. Denote by} $\underline{L}(S)$ \uline{the strictly full additive subcategory of} $\Qcoh(S)$ (2.2.1) \uline{generated by the} $L_i^{\otimes-n}$ \uline{for} $(i,n)\in I\times\mathbb N^+$\uline{. Let} $E\in\Ob\mathrm D(\Qcoh(S))$.
\begin{enumerate}[label=\alph*)]
\item \uline{The categories} $\underline{L}(S)$ \uline{and} $\Loclib(S)$ (2.0) \uline{are quasi-liftable in} $\Qcoh(S)$ (I 1.2) \uline{and quasi-stable under kernels of epimorphisms} (I 1.2).
\item \uline{Let} $n\in\mathbb Z$\uline{. For} $E$ \uline{to be} $n$\uline{-pseudo-coherent (resp.\ pseudo-coherent)} (2.2.1), \uline{it is necessary and sufficient that} $E$ \uline{be} $n$\uline{-pseudo-coherent (resp.\ pseudo-coherent) relative to} $\underline{L}(S)$ \uline{or} $\Loclib(S)$ (I 2.3)\uline{. The canonical functors}
\[
\mathrm K^{-}(\underline{L}(S))/\mathrm K^{-,\varnothing}(\underline{L}(S))\longrightarrow \mathrm D(\Qcoh(S))
\]
\uline{and}
\[
\mathrm K^{-}(\Loclib(S))/\mathrm K^{-,\varnothing}(\Loclib(S))\longrightarrow \mathrm D(\Qcoh(S))
\]
\uline{are fully faithful and have as essential image} $\mathrm D^{-}(\Qcoh(S))_{\coh}$ (2.2.1).
\item \uline{Let} $[m,n]$ \uline{be an interval in} $\mathbb Z$\uline{. For} $E$ \uline{to have perfect amplitude contained in} $[m,n]$ (2.2.1)\uline{, it is necessary and sufficient that} $E$ \uline{be isomorphic, in} $\mathrm D(\Qcoh(S))$\uline{, to a complex} $F$ \uline{such that} $F^i=0$ \uline{for} $i\notin[m,n]$, $F^i\in\Ob \underline{L}(S)$ \uline{for} $i>m$, \uline{with} $F^m$ \uline{locally free of finite type. The canonical functor}
\[
\mathrm K^b(\Loclib(S))/\mathrm K^{b,\varnothing}(\Loclib(S))\longrightarrow \mathrm D(\Qcoh(S))
\]
%%%%%%%%%%  scan idx 187  |  folio 174  %%%%%%%%%%
\SGAPrintPage{174}%
is
fully faithful and has as essential image $\mathrm D(\Qcoh(S))_{\parf}$ (2.2.1).
\end{enumerate}

\uline{Proof.} Let $F\to G$ be an epimorphism in $\Qcoh(S)$, where $G$ is of finite type. Since $F$ is an inductive limit of its quasi-coherent subsheaves of finite type (EGA I 9.4.9), there exists, by (2.1.1 a)), a quasi-coherent subsheaf $F'$ of $F$ of finite type such that the composite $F'\to F\to G$ is an epimorphism. By (2.2.3 (ii)), there exists an epimorphism $E\to F'$, with $E\in\Ob \underline{L}(S)$, so the composite $E\to F'\to F$, composed with $F\to G$, is an epimorphism, which proves that the category of quasi-coherent $\underline{O}_S$-Modules of finite type is quasi-liftable in $\Qcoh(S)$; a fortiori, the same holds for the category $\Loclib(S)$ or $\underline{L}(S)$. This proves the first assertion of a). For the second, we apply (1.1), taking for $\underline{S}$ the Zariski site of $S$, and for $\underline{C}$ (resp. $\underline{C}_0$) the fibered category $U\mapsto\Qcoh(U)$ (resp. $\underline{L}(U)$) ($U$ denotes a variable open subset of $S$, and $\underline{L}(U)$ the strictly full additive subcategory of $\Qcoh(U)$ generated by the $(L_i^{\otimes-n}\,|\,U)$). Since every object of $\underline{L}$ is pseudo-coherent, the notions of $n$-pseudo-coherence (resp. pseudo-coherence) (2.2.1) in the fibered category $\mathrm D(\Qcoh)$ coincide with the same notions relative to $\underline{L}$ (I 2.14). Condition (ii) (in the non-resp.\ case) of (1.1) being satisfied by (2.2.3 (ii)), we deduce, by (1.1 a)), the second part of a) concerning $\underline{L}(S)$; the assertion concerning $\Loclib(S)$ is proved similarly. Assertions b) follow immediately from (1.1 b) and c)), taking into account the fact that, when $E\in\Ob\mathrm D(\Qcoh(S))$ is $n$-pseudo-coherent, $E$ is an object of $\mathrm D^{-}(\Qcoh(S))$, since $S$ is quasi-compact. An object of the fibered category $\Qcoh$ is locally in $\underline{L}$ precisely when it is locally free of finite type: the category $\underline{L}$ is therefore locally liftable in $\Qcoh$ and locally stable under kernels of epimorphisms. Consequently, assertions c) follow from (1.2) applied with $\underline{C}=\Qcoh$, $\underline{C}_0=\underline{L}$, $\underline{C}_{1S}=\Loclib(S)$. This completes the proof.

\sgaStableTarget{sga6.e02.prop.2-2-9}\textbf{Proposition 2.2.9.} \uline{The hypotheses and notation are those of} (2.2.8)\uline{, but we further assume that} $S$ \uline{is separated or noetherian.}
\begin{enumerate}[label=\alph*)]
\item \uline{The canonical functors}
\[
\mathrm K^{-,b}(\underline{L}(S))/\mathrm K^{-,\varnothing}(\underline{L}(S))\longrightarrow\mathrm D(S)
\]
and
\[
\mathrm K^{-,b}(\Loclib(S))/\mathrm K^{-,\varnothing}(\Loclib(S))\longrightarrow\mathrm D(S)
\]
